% !PRINTED 701
\weberpage{701}
%<<CONT
und daher
\weq{20}{\mathop{\mathrm{Lim}}\limits_{s=1}\left(S - \frac{\gamma_s}{s-1}\right) = -\gamma\,\Gamma'(1) + D_1 = C_1,}
% [sic] The source prints gamma_s, +D_1, and =C_1. With the intended gamma^s, (19) gives the limit -gamma*Gamma'(1)-D_1. Definition (17), which subtracts gamma/(s-1), instead gives C_1=-gamma*Gamma'(1)+gamma*log(gamma)-D_1. The printed subscript, sign, and identification are retained and distinguished here.
wie zu beweisen war, endlich\footnote{Die hier gebrauchten Sätze über $\Gamma$-Functionen finden sich in den ausführlicheren Lehrbüchern der Integralrechnung, z.\,B. \gesperrt{Serret-Harnack}, Lehrbuch der Differential- und Integralrechnung, Bd.\ II, S.\,170\,f. (Leipzig 1885).}.

\begin{center}
\S.\,192.\\[4pt]
\gesperrt{Anwendung auf die Bestimmung der Classenzahl.}
\end{center}

Wir machen von diesen Sätzen jetzt die Anwendung auf die Theorie des Körpers $\Omega$. Wir verstehen unter den Zahlen $\mu_n$ des Theorems 5. die Normen der sämmtlichen Ideale des Körpers $\Omega$, also unter $T$ die Anzahl der Ideale in $\Omega$, deren Norm nicht grösser als eine gegebene positive Grösse $t$ ist, und erhalten
\weq{1}{\mathop{\mathrm{Lim}}\limits_{s=1}\sum \frac{s-1}{N(\mathfrak{a})^s} = \mathop{\mathrm{Lim}}\limits_{t=\infty}\frac{T}{t},}
und darin erstreckt sich die Summe der linken Seite auf \gesperrt{alle} Ideale $\mathfrak{a}$ des Körpers. Die Ideale zerfallen nun nach \S.\,153 in eine endliche Anzahl von Classen
\[
A_1, A_2, \ldots, A_h,
\]
und das grosse Ziel ist die Bestimmung dieser Zahl $h$, der \gesperrt{Classenzahl}.

Die Ideale $\mathfrak{a}_1$ einer dieser Classen $A_1$ sind dadurch charakterisirt, dass ihre Producte mit einem und demselben ganzen Functionale $\varphi_1$, das der Classe $A_1^{-1}$ angehört, \gesperrt{Hauptideale} sind, dass also
\[
\varphi_1 \mathfrak{a}_1 = \alpha
\]
eine ganze Zahl des Körpers $\Omega$ ist.

Ist die Norm von $\mathfrak{a}_1$ nicht grösser als $t$, so ist
\[
N_a(\alpha) \leqq N_a(\varphi_1)\, t = t_1.
\]
Ist $T_1$ die Anzahl der Ideale der Classe $A_1$, deren Normen nicht grösser als $t$ sind, so ist $T_1$ zugleich die Anzahl der nicht
%CONT>>
